\documentclass[10pt,a4paper]{amsart}
\usepackage[margin=19mm]{geometry}
\usepackage{amsmath,amssymb,mathtools}
\usepackage[scr=rsfs,cal=euler]{mathalfa}
\usepackage{xfrac}
\usepackage[unicode,hidelinks,bookmarks=false]{hyperref}
\newcommand{\ie}{i.e.\ }
\newcommand{\cf}{cf.\ }
\newcommand{\resp}{respectively\ }
\newcommand{\Subsection}{\subsection}
\providecommand{\nobreakdash}{\nobreak\hbox{-}}
\setlength{\emergencystretch}{2em}
\multlinegap0pt
\usepackage{xcolor}
\usepackage{csquotes}
\usepackage{amssymb}
\usepackage{mathtools}
\usepackage[shortlabels]{enumitem}
\usepackage{tikz}
\usepackage{float}
\newtheorem{prop}{Proposition}
\newcommand{\A}{\mathcal{A}}
\newcommand{\mA}{{\mathcal A}}
\title{Hamiltonian Cycles in Simplicial and Supersolvable Hyperplane Arrangements}
\author{Veronika K\"orber, Tobias Schnieders, Jan Stricker, Jasmin Walizadeh}
\date{Source: arXiv:2508.14538v3 (CC BY 4.0)}
\begin{document}
\maketitle

\noindent\emph{Source and licence.} Hamiltonian Cycles in Simplicial and Supersolvable Hyperplane Arrangements. Version arXiv:2508.14538v3 (CC BY 4.0), distributed by arXiv under the Creative Commons Attribution 4.0 International licence, http://creativecommons.org/licenses/by/4.0/, as recorded in the arXiv licence metadata for this exact version and shown on the abstract page https://arxiv.org/abs/2508.14538v3. Full licence notice: \texttt{licenses/arxiv-2508.14538-CC-BY-4.0.txt}.\par\smallskip

\noindent\textit{Editorial scope.} Counted unit: the article's prop prodHam (source0.tex lines 271--273). The article's complete original proof of it is delivered with it (source0.tex lines 274--285, one \texttt{proof} environment of 12 lines). No mathematics is written, repaired or re-derived: the statement and the proof are the article's own lines, in the article's own notation, and no retained line is substituted or re-flowed.  No further source-local context is needed: every cross-reference of every retained line resolves inside the two delivered spans.  Nothing was omitted from any retained span, so the package carries no omission record.  No retained line contains a citation, so the wrapper carries no bibliography and no bibliography entry.  No retained line contains a figure environment, an image-inclusion command or drawing code, so no image is delivered and the package declares no asset dependency.  Editorial formatting only, in addition: an AMS \texttt{amsart} preamble that restates the article's own declarations and macros used by the retained lines, namely the environments \texttt{prop} on separate counters, together with the standard packages those lines need.  The retained environments print in this document as Proposition 1 in order of appearance, whereas the article prints the same items with the numbering of its own full document.  The complete native source of the article as distributed in the arXiv e-print for this exact version is bundled unchanged as \texttt{sources/183558/source0.tex}.
	\begin{prop}\label{prodHam}
		If a hyperplane arrangement $\mA$ is the product of two hyperplane arrangements $\A_1$ and $\A_2$ that have a Hamiltonian cycle, then $\mA$ has a Hamiltonian cycle.
	\end{prop}

	\begin{proof}
		Suppose the hyperplane arrangement $\mA$ is the product of two hyperplane arrangements $\A_1$ and $\A_2$ with Hamiltonian cycles $\{a_0, a_1, \ldots , a_r \}$ and $\{b_0, b_1, \ldots , b_s \}$, respectively.
		Because the tope graph of $\mA$ is the product of the tope graphs of $\mA_1$ and $\mA_2$, the cycle
		\begin{align*}
			\{&(a_0, b_0), (a_0,b_1), \ldots,(a_0,b_s),\\
			&(a_1, b_s), (a_1,b_{s-1}), \ldots,(a_1,b_0),\\
			&(a_2, b_0), (a_2,b_1), \ldots,(a_2,b_s),\\
			&(a_3, b_s), (a_3,b_{s-1}), \ldots, \\
			&\ldots, (a_r,b_0)\}
		\end{align*}
		(reading from left to right) is a Hamiltonian cycle for $\A_1 \times \A_2$.
	\end{proof}

\end{document}
